% ID: SEL-D0037-OA-E04
% ESTADO: cotejada. Segunda pasada visual y borrador compilado revisados.
% Fuente: Selectividad/matematicas_aplicadas/Ua2MATEMÁTICAS APLICADAS examen.pdf
% SHA-256: a56739d297d4dc220d6898b24cf729dc556401a6bb292b8ca7cfb58c01c31651
% Página: 1; opción A; ejercicio original 4.

\textbf{(2.5 puntos)} Una característica poblacional $X$ sigue una distribución Normal $N(\mu,2.1)$. Sobre ella se formula un contraste de hipótesis bilateral con $H_0:\mu=5.5$ a un nivel de significación del 8\%. Se extrae una muestra aleatoria simple de tamaño 25 que proporciona una media muestral de 6.3. Plantee dicho contraste, determine su región crítica y razone si se puede aceptar la hipótesis nula.
